\scanpage{86}
% letterhead
UNIVERSITE DES SCIENCES \hfill Montpellier, le 17 Juillet 1975\\
ET TECHNIQUES DU LANGUEDOC\\
\texttt{-{}-{}-{}-{}-{}-{}-{}-{}-{}-}\\
Service de la Scolarité\\
AMC/GJ

\uline{Diffusion} : Liste H\\
N°4617

AVIS DE SOUTENANCE\\
D'UNE THESE DE DOCTORAT D'ETAT

\begin{center}\rule{2cm}{0.4pt}\end{center}

% printed form
\uline{Nom et prénom du candidat} : FOLLET Michel

\uline{Titre de la Thèse} : Etude physicochimique de la réactivité du cycle triazine : Synthèse et analyse spectrale de dérivés bicycliques.

\uline{Date, heure et lieu de soutenance prévus} :\\
8 Septembre 1975 à 17 Heures,\\
Salle de cours n° 9.

\uline{Nom du Directeur de thèse} : Professeur JACQUIER.

\uline{Désignation du laboratoire où la thèse est préparée} : Laboratoire de Synthèse et d'Etude Physicochimique d'Hétérocycles Azotés.

\uline{Résumé de la thèse} : Ce travail porte sur la série de la triazine - 1,2,4 , plus particulièrement sur les dérivés de l'acide thioxo-3 oxo-5 tétrahydro-2,3,4,5 triazine carboxylique -6. Différentes méthodes de synthèse de triazolo-triazines et triazolyl - triazines sont données. L'analyse spectrale de ces composés est faite par UV, RMN et spectrométrie de masse.

\uline{Durée approximative des recherches ayant abouti à la thèse} : 4 ans.

\uline{Noms des rapporteurs} :

JACQUIER R., Professeur, Chimie organique, U.S.T.L.\\
GRANGER R., Professeur, Chimie organique; Université MONTPELLIER I.\\
VIALLEFONT Ph. Professeur, Chimie organique, U.S.T.L.

\uline{Composition du jury} :

JACQUIER R., Professeur, Chimie organique, U.S.T.L. (Président)

GRANGER R., Professeur, Chimie organique, Université MONTPELLIER I (Membre extérieur choisi en raison de sa compétence scientifique)

DESCHAMPS J. Professeur, Chimie-physique, Faculté de Pau et des Pays de l'Adour (Membre étranger à la spécialité du candidat)

VIALLEFONT Ph. Professeur, Chimie organique, U.S.T.L.

DAUNIS J., Chargé de Recherche CNRS, Chimie organique, U.S.T.L.

LE SECRETAIRE GENERAL\\
\handwritten{\unsure{E S}}\\
\uline{E. SIAU}

% address label, lower left
MR LE PROFESSEUR GROTHENDIECK A.\\
MATHEMATIQUES \hfill 5-856
% archive sheet number: 85
